\chapter{Conclusions \& Future Work}
\label{chap:conclusions}

This chapter will conclude the dissertation, and discuss potential future work of the proposed Irish college course RS.

\section{Conclusions}

Choosing a suitable college course is a difficult, multi-faceted decision,
with 31\% of recent Irish college graduates reporting they were not likely to study their college course again, according to \cite{hea2018graduate}.
In the Irish context, GC's are turning their priorities from career counselling towards mental health counselling,
signalling the applicability of a technology in this domain.
In this light of this, we proposed an Irish college course RS to assist students in their college course exploration,
by providing them with diverse, relevant and serendipitous recommendations: all qualities of an effective RS.

Designing an effective college course RS requires a strong grasp on relevant domain knowledge.
This prompted a broad and comprehensive review of the fundamentals of college course choice,
in which we provided an answer to the question: \emph{``What college course should people study?''}.
Additionally, during our examination of closely-related works, we identified a potential research gap by unifying holland's RIASEC model with 
academic interests, as well as providing users with a smaller, more manageable recommendation set of specific college 
courses, as opposed to generalised areas of study.

Our review on the fundamentals of college course choice, along with interviews conducted with two professional GC's, 
were instrumental in our RS design.
This provided us with more clarity on the role of an RS in Irish guidance counselling,
desiring an RS that complements the human process of guidance counselling,
as opposed to replacing this process entirely.
This review also allowed us to identify the relevant information an RS should gather off a 
user in order to provide personalised college course recommendations.
This inspired the final design of this
dissertation's RS: a hybrid RS with Knowledge-based Filtering to filter college courses by NFQ levels, expected LC points 
and colleges, and a Content-based Filtering RS utilising academic interests and the RIASEC model
to provide 20 personalised college course recommendations.
Courses were proportionately recommended across the user's top-5 academic interests,
and courses with identical titles were not recommended.
All of which were efforts to enhance the diversity and serendipity of the recommendations.

This dissertation's RS was evaluated by conducting an experiment with $n=31$ participants.
An objective baseline RS was designed to return courses sorted descending by LC points, 
satisfying the user's explicit filters.
Compared to the baseline, the actual RS demonstrated significant improvements across the trust, serendipity and precision metrics,
with competitive performance on the diversity metric.
On average, the user liked 9 out of 20 college course recommendations provided by the actual RS,
where 8 of those recommendations were not on the user's original CAO list, i.e. serendipitous recommendations.
This is compared to the baseline's 4 out of 20 serendipitous recommendations,
indicating satisfactory performance of this dissertation's RS in the context of the research question.

\subsection{Contribution to the Literature}

To demonstrate our contribution to the literature, 
we will augment tables \ref{tab:research-gap-1} and \ref{tab:research-gap-2} introduced in section \ref{sec:research-gap}
to include this dissertation's RS, which will be referenced as
\emph{Irish College Course Recommender System} (ICCRS).

\begin{table}[H]
    \centering
\caption{Comparing features of all examined closely-related work (Part 1)}
\begin{tabular}{ |p{4cm}|p{3cm}|p{3cm}|p{5cm}|}
 \hline
 \textbf{Name} & \textbf{\# Questions} & \textbf{\# Results} & \textbf{Recommendation Type}\\
 \hline
\emph{Superprof}    & 10     & 3        & General area of study\\
\hline
\emph{Anderson Uni}     & 25     & 50       & College major\\
\hline
\emph{Marquette Uni}    & 10     & 30       & College major\\
\hline
\emph{\cite{qamhieh2020pcrs}}         & 125 & 7        & Engineering course\\
\hline
\emph{\cite{ghuge2023envisioning}}  & 6      & 5        & College major\\
\hline
\emph{\cite{lahoud2023comparative}}  & 6  & 5     & Lebanese college course\\
\hline
\emph{\cite{su2021personality}}  & 90     & 5  & Taiwanese college course\\
\hline
\emph{FMCC}         & $\sim$150 & 200   & Irish college course\\
\hline
\textbf{ICCRS}         & \textbf{126} & \textbf{20}   & \textbf{Irish college course}\\
\hline
\end{tabular}
\caption*{\textbf{Explanation of columns:}\\\textbf{\# Questions:} the number of questions asked to the user prior to receiving their college course recommendations.\\\textbf{\# Results:} The number of college course recommendations the user receives.\\\textbf{Recommendation Type:} the type of college course recommendation the user receives.}
\label{tab:research-gap-1-filled}
\end{table}

In observing table \ref{tab:research-gap-1-filled} above, ICCRS features a quiz with significant length,
demonstrating its breadth at 126 questions.
Additionally, inspired by the findings of \cite{toffler1970future} and \cite{long2025choice} in which smaller recommendation
sets yield more performant results, 20 college course recommendations are provided to the user.
To the best of our knowledge when considering the Irish context, providing the user with a more refined set of college course recommendations (20) 
is a novel contribution to the literature.

\begin{table}[H]
    \centering
\caption{Comparing features of all examined closely-related work (Part 2)}
\begin{tabular}{ |l|p{4cm}|p{7cm}|}
 \hline
 \textbf{Name} & \textbf{Can filter results by college?} & \textbf{Profiled User Information} \\
 \hline
\emph{Superprof}    & \ding{55}    & Academic interests\\
\hline
\emph{Anderson Uni}     & \ding{55}\tablefootnote{The only college majors recommended are ones offered from Anderson University.}     & Academic interests\\
\hline
\emph{Marquette Uni}    & \ding{55}\tablefootnote{The only college majors recommended are ones offered from Marquette University.}     & Academic interests\\
\hline
\emph{\cite{qamhieh2020pcrs}}         & \ding{55} & Hobbies, grades, MBTI\tablefootnote{Myer-Briggs Type Indicator (MBTI); a personality test developed in \cite{myers1962myers}.}\\
\hline
\emph{\cite{ghuge2023envisioning}}  & \ding{55}      & Hobbies, grades\\
\hline
\emph{\cite{lahoud2023comparative}}  & \ding{51}  & Hobbies, gender, academic interests\\
\hline
\emph{\cite{su2021personality}} & \ding{55}\tablefootnote{The only college majors recommended are ones offered from Cheng Shiu University, Taiwan.}     & RIASEC personality model\\
\hline
\emph{FMCC}         & \ding{55}\tablefootnote{You can filter your colleges by region, for example just the colleges in Dublin, Munster, etc.} & Academic interests\\
\hline
\textbf{ICCRS}         & \ding{51} & \textbf{Academic interests, RIASEC personality Model}\\
\hline
\end{tabular}
\caption*{\textbf{Explanation of columns:}\\\textbf{Can filter results by college?} is it possible for the user to specifically receive college course recommendations from colleges of their choice?\\\textbf{Profiled User Information:} What information is the system profiling off the user in order to make personalised college course recommendations?}
\label{tab:research-gap-2-filled}
\end{table}

In observing table \ref{tab:research-gap-2-filled} above,
ICCRS is one of the few works that allows users to filter college courses by particular colleges,
at least in the Irish context, ICCRS is the only RS capable of this.
To the best of our knowledge, utilising a user's RIASEC model and academic interests is a completely novel 
contribution to the literature surrounding college and career RS's.

\subsection{Revisiting the Research Question}

In section \ref{sec:research-question} we introduced this dissertation's research question:
\emph{``Can a Recommender System assist an Irish student in exploring college courses they would consider studying?"}.
Recall that we hypothesised that an RS that provides relevant, diverse and serendipitous recommendations would answer this research question.

In chapter \ref{chap:evaluation}, we compared the performance of an actual RS to a baseline. 
During our evaluation, we observed the actual RS slightly outperforming the baseline in terms of diversity,
which is satisfactory as the baseline RS is diverse by nature - as is evident by its 3.7/5 diversity score.
As well as this, we observed very large, statistically significant improvements 
for the precision and serendipity metrics.
On average, users responded that they would realistically consider studying 9 out of 20 college course recommendations,
8 of which are serendipitous recommendations not on the user's original CAO list.
This is compared to the baseline, where users only considered studying 4 out of 20 college recommendations,
all 4 of which were serendipitous recommendations,
indicating the actual RS was twice as effective in providing serendipitous recommendations compared to the baseline, which is a satisfactory result.

Following our evaluations from above, we can conclude that we have designed an RS that provides a user with relevant, diverse and serendipitous recommendations,
successfully answering this dissertation's research question.

\subsection{Revisiting Research Objectives}

In this subsection, we will revisit each research objective and examine whether the objective was achieved or not.

\begin{enumerate}
    \item \emph{Provide an answer to the question: ``What college course should people study?"}\\An answer to this question was provided in subsection \ref{sec:fundamentals-summary}, which largely inspired the design of this dissertation's RS. As a result, we can consider this research objective achieved.
    \item \emph{Identify the relevant information an RS should question a user to recommend suitable college courses.}\\The relevant information was identified in subsection \ref{sec:user-information-to-profile-summary}, directly informing the design of this dissertation's RS. Thus, we can also consider this research objective achieved.
    \item \emph{Design an objective, easily reproducable baseline RS appropriate to the Irish college context.}\\A baseline RS that fits this criteria was explicitly described in section \ref{sec:baseline-rs}. We may consider this research objective achieved.
    \item \emph{In comparison to a baseline college course RS, design an RS with statistically significant improvements for the user-perceived trust metric.}\\As observed in section \ref{sec:statistical-sig}, statistically significant improvements were found for the trust metric, when comparing the actual RS to the baseline. As a result, we can consider this research objective achieved.
    \item \emph{Analyse the qualitative insights and findings from a conducted experiment.}\\These insights were described in section \ref{sec:experiment-findings}, providing necessary qualitative analysis that goes beyond raw metrics. Therefore, we can consider this research objective achieved.
\end{enumerate}

With our evaluation of research objectives complete, we can now move on to discussing future works.

\section{Future Work}
\label{sec:future-work}

In this section, we will discuss a list of future works in the context of college course RS's.
Some of these future works will be relevant to the Irish context specifically, and
others will be relevant outside of the Irish context.

\subsection{Collaborative Filtering}

This dissertation's RS utilises Knowledge-based Filtering (KBF) and Content-based Filtering (CBF) in generating its recommendations.
According to \cite{aggarwal2016recommender}, CBF works by recommending items solely based on the user's provided preferences.
A disadvantage of this is what is known as a `filter bubble': 
users are only being recommended items with similar attributes to items they have already interacted with.
As a result, users find themselves in an `echo chamber', with recommendations lacking diversity and most importantly,
serendipity.
Another disadvantage of CBF is that the RS design requires significant metadata and labeling,
requiring a strong grasp on domain knowledge.
For instance in this dissertation's RS implementation, we required significant metadata regarding the academic interest
labels and RIASEC model associated with each college course.

On the other hand, Collaborative Filtering (CF) works by recommending items that other users with similar preferences also liked.
For example, recommending biology to a user because other users with similar preferences chose to study biology.
Due to the nature of its implementation, an advantage of CF lies in its ability to capture complex human preferences
that manually-labeled CBF would be unable to achieve.
For example, a CF RS may recommend a historical epic film to a user with past interest in sci-fi films,
as other sci-fi lovers enjoy long films on a grand scale - not only is this recommendation diverse, it may also be incredibly serendipitous.
This highly-nuanced information is incredibly difficult to encode with a dataset label,
posing a disadvantage for CBF.

Due to the unsupervised nature of CF, we hypothesise that a CF RS would be more optimal for a college 
course RS. Naturally, a limitation of the CF approach is the access to data regarding individual's preferences and their subsequent college course choice.
In chapter \ref{chap:initial-prototypes}, we explored a CF approach with a popular RIASEC dataset,
although the results were unsatisfactory due to a lack of coverage of quiz questions in the dataset.
To the best of our knowledge in the Irish context, this kind of data is unavailable,
posing a limitation to implementing CF in the Irish college course RS context.

\subsection{Additional Information to Profile off a User}

In section \ref{sec:fundamentals-summary}, we provided an answer to the question posed by research objective \ref{res-obj-answer}:
``\emph{What college course should people study?}".
The purpose of this was to provide insights into achieving research objective \ref{res-obj-information}:
``Identify the relevant information an RS should question a user to recommend suitable college courses''.
As part of our answer to research objective \ref{res-obj-answer},
we mentioned that the following two pieces of information would be relevant to profile off 
a user in the context of college course recommendations:

\begin{itemize}
    \item The user's desire for salary or employment after graduation.
    \item The user's desire for meaningful or socially impactful work.
\end{itemize}

In subsection \ref{sec:salary-design}, we discussed the viability of integrating salary or employment information into our college course RS.
Arguably, it would not be relevant to recommend a course with low mean graduate salaries to a student who 
highly prioritises employment rate and salary.
Given the lack of Irish course-specific data regarding salary or employment rates, this approach was not deemed possible.
That being said, as a piece of future work in the context of college course recommendations, we hypothesise that profiling a user's 
desire for salary or employment would result in more relevant recommendations.
In terms of effectively implementing this information, we hypothesise that a CF RS would be optimal for solving this problem,
due to the CF RS's ability to implicitly capture nuanced human preferences without explicit metadata.
For instance, if a user highly values salary and employment, with a CF RS they may be recommended college courses 
that other students who also highly valued salary and employment chose to study, such as medicine or law.
Due to the lack of Irish course-specific data for salaries and employment rates,
the CF approach is hypothetically applicable in this case.

Moreover, in subsection \ref{sec:meaning-design}, we discussed the viability of integrating a user's desire for meaning or social impact 
into our college course RS.
For instance, if an RS has knowledge that a student highly values social impact, the RS may be able to provide more tailored recommendations in light of this.
For example, the RS may recommend courses in social policy, nursing, or other fields with conventionally higher means for social impact.
That being said, this approach was turned down due to the difficulty in numerically translating or encoding this information as something more machine-readable.
However in terms of possible implementation ideas, we hypothesise that a CF RS may 
resolve this problem.
This is due to the CF RS's ability to capture nuanced human interaction that traditional metadata/labeling fail to capture.
For example, a CF RS may recommend a college course in social policy to a student due to their interest in meaningful work,
as other students who also valued meaningful work went into similar college courses.

\subsection{The Role of Large Language Models}

% future work - try without riasec? ngl academic interests does heavy lifting, lowkey curse of dimensionality, might not be necessary info/noise
    % but also talk about higher granularity of academic interest labels - reference the experiment insights here

% future work - allow the rec algo to run first, then ask the user to pick colleges
    % give example of people having a lack of awareness about colleges, especially specialised ones, e.g. marino, bimm, iadt, etc.
    % if they don't select those colleges, they won't get recs for it
    % * the RS could have a feedback loop, where the user is prompted again if they want to change their college
%     * e.g. declan only putting down tcd, ucd when he actually got recommended a lot of arts courses, could have benefitted from BIMM, IADT, etc.
%         * e.g. me wanting education but not doing dcu/marino/
%     * there are a looot of colleges and people are frankly unaware of what they do
%     * maybe the RS should run on all colleges (no filter, despite what the user puts down)- and if there are courses being recommended for those colleges, the system could ask the user: "hey, we got some course suggestions from these colleges because you scored highly in these areas, would you like to include these colleges in your final recommendation list?
%         * **when running test with doill he talked about this, mention it!!!**
%     * but it's a trade-off: if you over-engineer the system with feedback you can make the user drop out of the system - owen's advice

% implementation limitations
    % third language entry requirment for colleges
    % grade requirements for courses
    % inaccurate point representation for courses that require portfolios, tests or interviews





% \section{Large Language Models}

% maybe we should re-frame this - rather than taking the defensive angle and justifying the research, acknowledge it as future work.
% the brutal reality is that if utilised effectively, LLM's would blow my product out of the water, you simply can't compete.
% let's be honest here and dont be delusional

% owen - it's a black-box
    % can give you post-hoc 
    % human-verfiied cao database
        % llm's can hallucinate this info, e.g. recommend course that doesn't exist 

% The usage of Large Language Models (LLM's) has become ubiquitous in today's society.
% For instance in February 2026, 
% OpenAI reported 900 million\footnote{According to \url{https://openai.com/index/scaling-ai-for-everyone/}} 
% weekly active users on their LLM-powered chatbot \emph{ChatGPT}\footnote{https://chatgpt.com/}.
% According to \cite{yankouskaya2026lets}, over 40\% of adults in the UK are happy to use ChatGPT as a mental health counsellor,
% indicating a significant amount of trust in LLM's towards important life decisions.
% - most notably, choosing a college course is also an important life decision.
% Given the widespread use and selective, immense trust of LLM chatbots, 
% this section will answer the question: \emph{``Can a user just use an LLM to get college course recommendations? Is this research really necessary?"}

% As described in \cite{alansari2025large}, LLM's have the capability of hallucinating, or
% presenting false or fabricated information as facts.
% In the context of Irish college courses recommendations,
% this presents the possibility of an LLM providing inaccurate or 
% false information regarding Irish college courses - for example, the LLM may suggest a course that does not exist, or 
% provide inaccurate entry requirements.
% This calls for an RS with accurate, up-to-date information,
% ensuring that the college course recommendations are grounded in this information,
% preventing the possibility of hallucination.

% As we discovered in subsection \ref{sec:fundamentals},
% answering the question: \emph{``What college course should I study?"}
% is not straightforward at all.
% Generally, these LLM chatbots are designed to be open-ended with a simplified interface,
% putting more agency into the user's hands.